\section{Обсуждение}\label{sec:discussion}

\subsection{Сводка результатов}

Собранные данные позволяют дать следующие ответы на вопросы, поставленные
в разделе~\ref{sec:problem}.

\emph{Сохраняется ли фрактальность при добавлении инерции?} Да, полностью.
Размерность границы в сечении $(\eta,\beta)$ равна
$\DEtaBeta \pm \DEtaBetaS$ и совпадает в пределах погрешности с базовой картой
$(\eta_0,\eta_1)$, где $\DBaseTanh \pm \DBaseTanhS$. Инерция сдвигает границу,
меняет соотношение долей режимов, но не делает её гладкой.

\emph{Влияние затухания весов и размера батча.} Умеренное затухание весов
сглаживает границу; уменьшение размера батча, наоборот, увеличивает её
изрезанность монотонно ($\DMbFull \to \DMbFour \to \DMbOne$ при
$B = 32 \to 4 \to 1$). Оба эффекта естественно интерпретируются через
многомасштабность целевой функции: регуляризация подавляет мелкомасштабные
особенности ландшафта, стохастичность батча их добавляет.

\emph{Сглаживают ли адаптивные методы границу?} Ответ оказался тоньше, чем
предполагала исходная гипотеза, и составляет, на наш взгляд, основной
содержательный результат работы.

\subsection{О «сглаживании» границы адаптивными методами}

Гипотеза задания состояла в том, что адаптивные методы частично сглаживают
фрактальную структуру за счёт нормализации градиентов по координатам. Данные
подтверждают её в узком смысле и опровергают в широком.

В плоскости $(\eta_0,\eta_1)$ Adam действительно даёт меньшую размерность, чем
SGD ($\DAdam$ против $\DSgd$). Однако наименьшую размерность даёт не адаптивный
метод, а SGD с фиксированной инерцией $\beta=0.9$ ($\DMomentum$), где никакой
покоординатной нормализации нет. Значит, сглаживание обеспечивает не
адаптивность как таковая, а усреднение градиента по времени, которое подавляет
высокочастотную составляющую динамики и уменьшает чувствительность исхода
к малым изменениям шага.

Более того, сглаживание оказывается свойством выбранного \emph{сечения}, а не
оптимизатора. Стоит перерезать пространство гиперпараметров вдоль собственных
осей Adam, и граница вновь оказывается предельно изрезанной: $\DAdamBOne$ для
$(\eta,\beta_1)$ и $\DAdamBTwo$ для $(\eta,\beta_2)$ --- выше, чем у SGD в
базовом сечении. Механизм понятен из вида карт Adam и RMSprop
(рис.~\ref{fig:opt}): покоординатная нормализация делает величину шага по
каждому весу почти независимой от масштаба градиента, из-за чего динамика
разделяется по слоям и карта в плоскости $(\eta_0,\eta_1)$ становится почти
сепарабельной --- граница вырождается в набор почти горизонтальных полос.
Низкая размерность здесь измеряет сепарабельность, а не отсутствие хаоса.
Как только сечение выбирается так, что сепарабельность не работает
(оси $\beta_1$, $\beta_2$ управляют самим усреднением), фрактальность
проявляется в полной мере.

Этот вывод согласуется с результатами~\cite{torkamandi2025}, где фрактальные
границы наблюдались именно для моделей, обучаемых Adam, и
с~\cite{cohen2022}, где показано, что адаптивные методы обладают собственным
аналогом грани устойчивости. Практический смысл: переход на Adam не избавляет
от необходимости аккуратного подбора шага обучения, он лишь переносит
нерегулярность в другие координаты пространства гиперпараметров.

\subsection{Связь с теорией динамики}

Наблюдаемая картина хорошо укладывается в схему раздела~\ref{sec:theory}.
Для квадратичной модели граница устойчивости --- гладкая прямая
$\eta\lambda = 2(1+\beta)$. В карте $(\eta,\beta)$ (рис.~\ref{fig:etabeta})
внешняя граница области переполнения действительно близка к прямой линии,
смещающейся с ростом $\beta$; именно она соответствует линейной теории.
Вся нерегулярность сосредоточена на внутренней границе --- между сходимостью
и ограниченными колебаниями. Это в точности та зона, где линейный анализ
неприменим: система вышла за порог устойчивости по старшей моде, но не
расходится, потому что нелинейность переводит её в область меньшей
кривизны (catapult-механизм~\cite{lewkowycz2020}) или удерживает на грани
устойчивости~\cite{cohen2021}. Исход в этой зоне определяется тонким балансом
и потому чувствителен к сколь угодно малым изменениям гиперпараметров ---
что и порождает фрактальную структуру.

Бифуркационная диаграмма (рис.~\ref{fig:bif}) показывает механизм в чистом
виде: при росте $\eta$ нелинейная скалярная система проходит через каскад
удвоения периода и хаос, и рост $\beta$ смещает каскад в область меньших
шагов. Показательно, что для линейной задачи инерция допустимый шаг
расширяет (следствие 2), а для нелинейной --- сужает окно регулярного
поведения. Именно это расхождение между линейным предсказанием и нелинейной
реальностью и есть источник изрезанности границы.

\subsection{Ограничения исследования}\label{sec:limitations}

Результаты следует интерпретировать с учётом ряда существенных ограничений.

\textbf{Зависимость оценки размерности от разрешения сетки.} Это наиболее
серьёзное методологическое ограничение. Одна и та же конфигурация, посчитанная
на разных сетках, даёт систематически разные оценки (табл.~\ref{tab:res}):
от $\DResSmall$ на сетке $128\times128$ до $\DBaseTanh$ на $512\times512$.
Причина в том, что более мелкая сетка разрешает более тонкие складки границы,
а диапазон доступных масштабов $\varepsilon$ в формуле~\eqref{eq:box} всегда
ограничен снизу размером пикселя. Приводимая величина $\sigma_D$ отражает
только качество линейной регрессии и эту систематику не учитывает. Поэтому
сравнивать между собой корректно только карты одинакового разрешения, что
в настоящей работе и делается: все сравнения оптимизаторов, сечений и
датасетов выполнены на сетках $256\times256$.

\begin{table}[htbp]
  \centering
  \caption{Зависимость оценки размерности от разрешения сетки для одной и той
  же конфигурации (tanh, $(\eta_0,\eta_1)$, $|\mathcal{D}|=1$).}
  \label{tab:res}
  \small
  \input{tables/tab_resolution.tex}
\end{table}

\textbf{Различие окон сканирования.} Порог устойчивости зависит от активации,
объёма данных и оптимизатора, поэтому окна по $\eta$ подбирались
индивидуально. Как следствие, доля площади карты, занятая границей, различна
для разных экспериментов, а измеряемая размерность от неё зависит. Сравнение
размерностей между оптимизаторами является, строго говоря, качественным;
количественно надёжны лишь сравнения внутри одного семейства окон
(например, ряд по размеру батча, где окно фиксировано).

\textbf{Граница как зона, а не кривая.} В отличие от классических фракталов,
где граница --- одномерная кривая, в наших картах переход между сходимостью и
колебаниями представляет собой протяжённую зону, в которой режимы перемешаны
на масштабе отдельных узлов сетки. Значения $D\approx1.85$, полученные на
глубоких уровнях зума, отражают приближение этой зоны к площадному заполнению,
а не структуру гладкой кривой. Формально box-counting применим и здесь, но
интерпретировать результат как «размерность границы» следует с осторожностью.

\textbf{Артефакт критерия сходимости для ReLU.} Критерий, унаследованный от
предыдущей НИР (средний нормализованный loss меньше единицы), засчитывает как
успешное обучение вырожденный случай: при большом шаге все ReLU-нейроны
переходят в область отрицательных предактиваций, градиент обнуляется, обучение
останавливается, и loss замирает на постоянном уровне ниже начального. Мёртвая
сеть формально классифицируется как сошедшаяся. Именно этим объясняется доля
сходимости $0.99$ для ReLU при $|\mathcal{D}|=1$ во всём диапазоне вплоть до
$\eta = 10^{4}$. Для корректного сравнения активаций пришлось использовать
набор из 32 примеров, где такой коллапс менее вероятен. Более строгий критерий
(требование снижения loss в заданное число раз) устранил бы артефакт, но
сделал бы результаты несопоставимыми с предыдущей работой.

\textbf{Ограниченность числа итераций.} Все эксперименты выполнены при
$T = 400$. Классификация режимов зависит от $T$: траектория, помеченная как
осциллирующая, могла бы при большем $T$ выйти на сходимость. Это смещает
границу в сторону области колебаний, но, судя по устойчивости картины при
переходе от $T=300$ к $T=500$ на этапе калибровки, не меняет качественных
выводов.

\textbf{Простота архитектуры и объём данных.} Исследована однослойная сеть с
512 параметрами на выборках объёмом до 32 примеров. Перенос выводов на
глубокие архитектуры и полноразмерные датасеты требует отдельной проверки;
косвенным аргументом в пользу переносимости служит работа~\cite{torkamandi2025},
где явление наблюдалось для трансформеров.

\textbf{Отсутствие прямого трёхмерного box-counting.} Трёхмерное пространство
исследовано набором двумерных сечений; как показано в разделе~\ref{sec:3d},
немонотонность $D(\lambda_{wd})$ не позволяет по сечениям надёжно судить о
размерности трёхмерной границы.

\subsection{Практические следствия}

Из полученных данных вытекает несколько практических соображений.
Во-первых, автоматический подбор шага обучения принципиально осложнён тем, что
на любом масштабе рядом с работающей конфигурацией находится неработающая;
это относится и к SGD, и к Adam, только для последнего нерегулярность
проявляется в других координатах. Во-вторых, наилучшие по скорости
конфигурации систематически лежат вблизи границы, что согласуется с концепцией
грани устойчивости и означает, что консервативный выбор шага заведомо жертвует
скоростью. В-третьих, умеренная $L_2$-регуляризация и увеличение размера батча
делают ландшафт обучаемости более предсказуемым --- это может быть
самостоятельным аргументом в их пользу помимо влияния на обобщающую
способность.
